\documentclass[11pt]{article}
\usepackage[margin=1in]{geometry}
\usepackage{authblk}
\usepackage{graphicx}
\usepackage{booktabs}
\usepackage{amsmath, amssymb}
\usepackage{caption}
\captionsetup{font=small,labelfont=bf}
\usepackage{array}
\usepackage[hidelinks]{hyperref}

\title{$\Delta KL$ as an innovation measure for trademark filings:\\
the expected reward to forward flux, and the unexpected reward to
backward flux}
\author{Eric Silver}
\affil{}
\date{May 2026}

\begin{document}
\maketitle

\begin{abstract}
We compute a Kullback--Leibler surprise difference $\Delta KL =
KL_{\text{pros}} - KL_{\text{retr}}$ on each of $10{,}060{,}073$
U.S.\ trademark filings in $1990$--$2021$ against per-day
exponentially decay-weighted references of the same NICE class
(half-life $H = 2$ years; chosen by signal-maximising sweep).
Trademark filings are a new substrate for the
prospective/retrospective surprise framework of Murdock, Allen, and
DeDeo (2017): they are fee-bearing, legally consequential, and
filed under economic incentive, so the goods/services prose
reflects the firm's deliberate language choice rather than
incidental editorial drift. $\Delta KL$ is essentially uncorrelated
with the dominant alternative innovation proxy --- firm-year patent
counts, $r = +0.010$ on $174{,}569$ firm-years --- and the magnitude
of its correlation with $4$-year stock excess returns is comparable
to (slightly larger than) that of patents: incremental $R^2$ over
controls of $+0.0046$ for $\Delta KL$ vs.\ $+0.0024$ for $\log(1 +
n_\text{patents})$. Including both predictors yields $R^2$ gain of
$+0.0069$, essentially the sum of the two separate contributions:
the two innovation channels are independent and of comparable
magnitude.

\textbf{Where $\Delta KL$ aligns with the distinctiveness literature.}
On debut filings, the response curve of $5$-year post-registration
survival on the innovation side ($\Delta KL > 0$) is a clean
saturating exponential: $P(\text{survived}) = 0.995 - 0.55\,
e^{-0.063\,\Delta KL}$, $R^2 = 0.499$ across $20$ binned points.
Survival rises with prospective novelty, concave-down, with
diminishing returns. This is the empirical content of optimal
distinctiveness (Deephouse, 1999; Zhao, Fisher, Lounsbury, \&
Miller, 2017) translated from \emph{level} of category-typicality
to \emph{directional} vocabulary flux.

\textbf{Where it diverges.} The strict inverted-$U$ prediction of
the optimal-distinctiveness literature ($U$-shaped pay-off on
distinctiveness level, with a turning point at high distinctiveness)
does not appear: the response curve keeps rising at the positive
extreme, with no inflection. More surprisingly, the harvest side
($\Delta KL < 0$) has its own saturating exponential, fitted much
more tightly ($R^2 = 0.978$): filings whose vocabulary lies further
into the receding/legacy side of the category are not penalised but
rewarded, and the reward grows the deeper the harvest. We read this
as \emph{increasing returns to late-stage imitation}, not predicted
by the distinctiveness literature. Together the two saturating
exponentials produce a survival $U$ on $\Delta KL$ in the pooled
corpus; the $+U$ holds in $32$ of $44$ NICE classes, only one
industry has a significant inverted-$U$, and $11$ are indeterminate.
We rule out nine alternative explanations and identify a
cluster-viability mediator that absorbs $84\%$ of the quadratic on
survival, leaving a residual linear $\Delta KL$ effect intact.

\textbf{An unresolved asymmetry: harvest dominates the
public-listing signal.} The probability that a filing's owner ever
appears in SEC EDGAR (proxy for IPO/public listing) is, at the
filing level, monotone \emph{decreasing} in $\Delta KL$ rather than
$U$-shaped: $5.78\%$ at the most-negative $\Delta KL$ decile, $4.0\%$
at the flux-neutral middle, $3.49\%$ at the most-positive decile.
The harvest tail captures the EDGAR-listing pattern far more cleanly
than the innovation tail does. We do not have a single-mechanism
explanation for why late-stage-imitation filings are so much more
strongly associated with the public-listing outcome than novelty
filings are at the unit of one filing.

A long companion paper (\texttt{main.pdf}) reports the per-industry
detail, the financial-outcomes panel on $\sim 1{,}500$ SEC-listed
firms, the patent-counts crosswalk, the cluster-viability mediation,
and the per-NICE-class deep dives.
\end{abstract}

\section{$\Delta KL$: construction, novelty of the substrate, and relation to existing innovation measures}

\paragraph{Construction.} For each filing $i$ at filing-date $t_i$
in NICE class $c$, the goods/services term distribution $P_i$ is
compared via support-restricted KL divergence to two reference
distributions $Q_i^{\text{pros}}$ and $Q_i^{\text{retr}}$ built by
per-day exponential decay-weighting (half-life $H = 2$ years; the
signal-maximising choice across a sweep of $H \in \{0.5,1,2,3,5\}$
years on Class~9) of all other filings in $c$ before and after
$t_i$ respectively. $\Delta KL_i = KL_i^{\text{pros}} -
KL_i^{\text{retr}}$ is positive when filing $i$'s vocabulary is
unusual against the recent past \emph{and} closer to the near
future --- the directional resonance signal of Murdock, Allen, and
DeDeo (2017) applied to commercially incentivised text. The two
component levels are correlated at $\rho = 0.99$ per filing; the
difference is the directional residual.

\paragraph{Why trademarks.} Trademark goods/services prose is a new
substrate for the prospective/retrospective surprise framework, and
its properties make it well-suited to the measurement question.
Filings are \emph{fee-bearing} (each U.S.\ class costs the applicant
\$250-\$350; abandonment penalties accrue), \emph{legally
consequential} (granted marks become enforceable property rights;
non-use within five years is grounds for cancellation), and
\emph{filed under economic incentive} (firms file when they intend
to commercialise, not when they're being literary). The prose
therefore reflects deliberate language choice by a commercial actor
with real money at stake on the front end and real legal exposure on
the back end. Prior applications of the surprise framework (Murdock,
Allen, \& DeDeo, 2017, on long-form essays; Barron, Huang, Spang, \&
DeDeo, 2018, on parliamentary speech) have used corpora where the
authors faced reputational but not direct economic cost for their
linguistic choices. The trademark corpus is the first such
application to incentivised commercial text at this scale
($10$+ million filings).

\paragraph{Relation to patents.} The dominant innovation measure in
the firm-level empirical literature is the patent count (Hall, Jaffe,
\& Trajtenberg, 2005; Kogan, Papanikolaou, Seru, \& Stoffman, 2017).
$\Delta KL$ is essentially uncorrelated with firm-year patent counts
($r = +0.010$ on $174{,}569$ firm-years), and the magnitude of its
correlation with $4$-year excess stock returns is comparable to (and
slightly larger than) that of patents. In the firm-year R$^2$
decomposition for predicting $4$-year excess returns over the S\&P
500: controls-only $R^2 = 0.0211$; adding $\log(1 + n_\text{patents})$
gives $R^2 = 0.0235$ (incremental $+0.0024$); adding $\Delta KL$
alone gives $R^2 = 0.0257$ (incremental $+0.0046$, roughly twice
patents); including both gives $R^2 = 0.0280$ (incremental $+0.0069$
over controls, essentially the sum of the two contributions). The
two innovation channels are independent and of comparable magnitude.
Patents proxy a firm's continuing R\&D output; $\Delta KL$ on the
firm's debut filing proxies whether the firm started in fertile
linguistic ground.

\paragraph{Relation to the distinctiveness literature, in advance.}
The closest-adjacent management-theory framework is optimal
distinctiveness (Deephouse, 1999; Zhao, Fisher, Lounsbury, \&
Miller, 2017), which predicts an inverted-$U$ payoff on the
\emph{level} of category-distinctiveness: too typical is invisible,
too distinctive is illegitimate. $\Delta KL$ is not a level measure
of distinctiveness; it is a \emph{directional} measure of where a
filing sits between recent past and near future of the category's
vocabulary. The results that follow align with the empirical content
of optimal distinctiveness on the positive-$\Delta KL$ side
(saturating returns to novelty) but diverge from its strict
inverted-$U$ prediction (no turn-down at the extreme) and add a
second saturating-exponential response curve on the harvest side
that the literature does not predict. We present both as results.

\section{Two results: optimal distinctiveness on the innovation tail; increasing rewards to late-stage imitation on the harvest tail}

\begin{figure}[t]
\centering
\includegraphics[width=\linewidth]{results/outcome_by_kl_lines_v2.png}
\caption{Three outcomes by KL bin (Wilson 95\% CIs as error bars),
pooled across all $44$ NICE classes; $n = 10{,}060{,}073$ clean
H=2y filings, $1990$--$2021$ cohorts. Each panel shows a different
shape on $\Delta KL$ (blue). \textbf{Panel A
($P(\text{reached registration})$):} $U$-shaped on $\Delta KL$,
asymmetric, with the harvest extreme ($54.8\%$) higher than the
innovation extreme ($47.5\%$) and the flux-neutral middle lowest
($\sim 45\%$). \textbf{Panel B ($P(\text{survived 5y} \mid
\text{reached registration})$):} \emph{monotone increasing} in
$\Delta KL$ --- once past the examiner, more prospective novelty
means more survival. The substantive post-registration mortality
signal favours the innovation side. \textbf{Panel C ($P(\text{owner
ever in SEC EDGAR})$):} \emph{monotone decreasing} in $\Delta KL$
($5.78\%$ at the most-negative decile to $3.49\%$ at the
most-positive decile). The public-listing signal is dominated by
the harvest tail; we discuss this unresolved asymmetry below. On
the component KLs (red, green), registration is monotone-up
(more-distinctive marks collide less with prior art and clear
examination more easily under $15$~U.S.C.\ \S\,1052(d)).}
\label{fig:short-headline}
\end{figure}

Figure~\ref{fig:short-headline} is the central empirical object.
The $U$ on \emph{joint} 5-year survival (reached registration
\emph{and} still live; not separately plotted but obtained by
multiplying Panel A by Panel B per bin) places the flux-neutral
middle of $\Delta KL$ lowest in the corpus at $37\%$ vs.\ $44\%$ at
the harvest extreme and $41\%$ at the innovation extreme. Both
extremes outperform the middle on this joint outcome; the
functional form of each tail tells a different theoretical story,
to which we now turn.

\paragraph{Innovation tail: optimal distinctiveness as a saturating
exponential.} On the debut-filings-only panel ($n = 5{,}088{,}700$
owner-first filings), we fit a saturating exponential of survival on
$|\Delta KL|$ on each side: $P(\text{survived}) = a +
b\,e^{-c|\Delta KL|}$ on $20$ binned points. On the positive side
($\Delta KL > 0$) the fit is $a = 0.995$, $b = -0.55$, $c = 0.063$
with $R^2 = 0.499$ across $20$ bins. The curve is monotone increasing
in $\Delta KL$, concave-down, saturating: an additional unit of
prospective-novelty pays at a decreasing rate. This is the empirical
content of optimal distinctiveness (Deephouse 1999, \emph{SMJ}; Zhao,
Fisher, Lounsbury \& Miller 2017, \emph{SMJ}) translated from
\emph{level} of distinctiveness to \emph{directional} vocabulary
flux: novelty pays with diminishing returns, no inflection where
``too distinctive'' becomes a penalty. The strict inverted-$U$
prediction of the original theory is not what we find, but the
underlying empirical claim --- that distinctiveness has positive but
saturating returns --- is supported, and the exponential is the
explicit functional form.

\paragraph{Harvest tail: increasing rewards to late-stage imitation.}
On the negative side ($\Delta KL < 0$) the same functional form fits
much more tightly: $a = 0.573$, $b = -0.137$, $c = 3.33$, $R^2 =
0.978$ across $20$ bins. Survival rises from $\sim 0.44$ at $\Delta
KL \approx 0$ to $\sim 0.57$ at the most extreme negative $\Delta
KL$. Conditional on negative $\Delta KL$ the Pearson correlation
between $|\Delta KL|$ and survival is $+0.054$ ($n = 2{,}586{,}126$
debuts; Spearman $+0.058$). The harvest tail is doing real work and
is unmistakably present. The substantive interpretation: filings
whose vocabulary lies further into the receding/legacy side of the
category are not penalised but rewarded, and the reward \emph{grows}
the deeper one harvests. The further past the category's modal
present-day vocabulary a debut filing sits, the more durable its
mark turns out to be. This is the late-stage-imitation analogue of
optimal distinctiveness: rather than diminishing returns to novelty,
the data show \emph{increasing} returns to picking up legacy
vocabulary as the field collectively abandons it. The asymmetric
saturation rates ($c = 0.063$ on the innovation side vs.\ $c = 3.33$
on the harvest side) say the harvest side approaches its asymptote
much faster, which is consistent with the small-firm-can-pick-up-
the-legacy-niche story: there are not that many firms staying in
legacy vocabulary; the ones who do are competing against few peers
and the marginal return to going even deeper into legacy is high
but quickly capped by how much of the niche exists at all.

\paragraph{The most striking single-token examples of the harvest
tail are in retreating industries.} Negative-$\Delta KL$ tokens that
survive at $50$--$80\%$ in the corpus: \texttt{microfiche},
\texttt{dial-up}, \texttt{pager}, \texttt{store-and-forward},
\texttt{lager}, \texttt{stout}, \texttt{shandy}, \texttt{annuity},
\texttt{microcomputers}. None of these are growing categories; all
support trademarks that the filing firm keeps alive long after the
vocabulary itself recedes from the corpus.

\section{How durable is the $U$? Industry $\times$ year breakdown}

Table~\ref{tab:short-classification} reports the $U$-shape
classification of every $($industry, year$)$ cell in the panel.
For each of $44$ industries with $\geq 5{,}000$ clean filings, we
fit a logistic regression of $5$y survival on $z(\Delta KL)$,
$z(\Delta KL)^2$, and $\log(n_{\text{terms}})$ separately for each
filing-year $1990$--$2021$, and classify the cell by whether the
quadratic coefficient is significantly positive (the $U$), significantly
negative (inverted-$U$), or indeterminate.

\begin{table}[h]
\centering
\small
\begin{tabular}{lrr}
\toprule
Statistic & count & \% of total \\
\midrule
\textbf{Industries (of 44 with sufficient data)} & & \\
  with significant $+U$ on $\Delta KL$ & 32 & 73\% \\
  with significant inverted-$U$ on $\Delta KL$ & 1 & 2\% \\
  indeterminate & 11 & 25\% \\
\midrule
\textbf{Industry $\times$ year cells (of 1,267 total)} & & \\
  significant $+U$ on $\Delta KL$ ($z > +1.96$) & 338 & 27\% \\
  significant inverted-$U$ ($z < -1.96$) & 81 & 6\% \\
  indeterminate quadratic & 848 & 67\% \\
  linear advantage $+\Delta KL$ ($z > +1.96$) & 197 & 16\% \\
  linear advantage $-\Delta KL$ ($z < -1.96$) & 484 & 38\% \\
  neither (linear $|z| \leq 1.96$) & 586 & 46\% \\
\bottomrule
\end{tabular}
\caption{$U$-shape classification of $44$ NICE classes individually
and of the $1{,}267$ $(\text{industry}, \text{year})$ cells where a
logit fit converged. The pooled $+U$ holds in $32$ of $44$ industries;
only $1$ has a significant inverted-$U$ (Telecommunications, where
the linear $\Delta KL$ effect is positive but the harvest tail is
absent). The remaining $11$ are not inverted-$U$ industries --- they
are indeterminate (no significant curvature in either direction).
Cell-level: $27\%$ of cells have a significant $+U$ on $\Delta KL$;
$6\%$ have a significant inverted-$U$; $67\%$ are indeterminate.
Linear-direction counting at the cell level shows a striking
asymmetry: $38\%$ of cells have a significant linear advantage to
\emph{negative} $\Delta KL$ (the harvest direction), versus $16\%$
with significant advantage to positive $\Delta KL$.}
\label{tab:short-classification}
\end{table}

\paragraph{Where does the $U$ not hold and why?} The single
significant inverted-$U$ industry is Telecommunications (NICE 038,
$n = 170{,}826$, $t_{\text{quad}} = -11.8$). It has a positive
linear $\Delta KL$ coefficient ($+0.015$): forward-flux still helps
in telecoms, but the harvest tail is absent --- old telecom
vocabulary does not survive in the way old typesetting or
old-beverage vocabulary does. The $11$ indeterminate industries
split: some (Scientific \& Tech Services $042$, Transport $039$)
have a positive linear coefficient on $\Delta KL$ but neither a
significant quadratic nor a strong harvest tail; others (Hand Tools,
Trim, Jewelry, Vehicles, Ropes \& Textiles, Musical Instruments)
have a \emph{negative} linear coefficient --- mild harvest advantage,
no significant innovation tail to complete the $U$. The interpretation
in both cases is that the industry advantages flux in one direction
but not the other, so the quadratic does not register as
significantly $U$-shaped. None of the $11$ are inverted-$U$
industries where firms in the flux-neutral middle outperform both
extremes.

\section{Nine alternative explanations ruled out; cluster-viability mediates}

Table~\ref{tab:short-falsification} reports the falsification battery
in condensed form. The headline result is that nine alternative
hypotheses for the $U$ are ruled out and a tenth check
(competition-density mediation) identifies the mediator that absorbs
$84\%$ of the quadratic on survival, leaving the linear $\Delta KL$
effect intact.

\begin{table}[h]
\centering
\small
\begin{tabular}{p{4.0cm}p{11.2cm}}
\toprule
\textbf{Alternative} & \textbf{Outcome} \\
\midrule
(1) Foreign-extension artefact (Madrid-Protocol marks) &
Rejected. US-only and foreign-only sub-samples give nearly identical
quadratics ($+0.0078$ and $+0.0100$). \\
\addlinespace
(2) First-time vs.\ recurring filings &
Both have the $U$; recurring slightly stronger. \\
\addlinespace
(3) Same $U$ on financial outcomes &
No. Firm-year gross margin is roughly linear in $\Delta KL$ ($+3.2$ pp/$\sigma$);
the U lives in survival. \\
\addlinespace
(4) Large-incumbent survivor bias on the harvest tail &
Partial. The bottom $\Delta KL$ decile has larger mean firm
filing-volumes than the middle. \\
\addlinespace
(5) ``Liquid Death'' hypothesis (low-$\Delta KL$ is novel products in old vocabulary) &
Rejected on text-length splits. \\
\addlinespace
(6) The $U$ is firm-level (firms with flux-neutral filings just fail) &
Rejected. Within firms that abandoned some marks and kept others
alive (the ``mixed'' cohort, $349{,}819$ firms), the $U$ on
filing-level survival is preserved. \\
\addlinespace
(7) Blockchain / NFT wave as imitated-but-failed counterexample &
Rejected. The crypto cohort outperforms the corpus on every survival
margin. \\
\addlinespace
(8) Vintage pooling artefact &
Rejected. Per-year fits 1990--2021 show the pooled quadratic is
significantly positive in $23$ of $32$ years independently; the $U$
does not track the macro cycle. \\
\addlinespace
(9) English-language usage drift (Google Books n-grams) &
Partial shared signal, not a confound. Per-token correlation between
trademark log-lift and Google Books English log-lift is $+0.47$
($n = 53$). About $22\%$ of trademark-flux variance per token is
shared with general English usage; $78\%$ is corpus-specific. \\
\midrule
\textbf{(10) Competition-density mediation} (the explanation that
sticks) &
$\log(1 + \text{live\_competitors})$ in an asymmetric
$[t-2y,\, t+5y]$ window (cosine $\geq 0.7$ on g/s TF-IDF; not the
H=2y reference window) absorbs $84\%$ of the $\Delta KL$ quadratic
on survival. Flux extremes are in vocabulary clusters where peers
also survive, not in less-competitive clusters. Linear $z(\Delta KL)$
effect remains positive ($z = +2.65$) after the mediator is
included; flux is partly but not wholly redundant with cluster
viability. \\
\bottomrule
\end{tabular}
\caption{Falsification battery for the survival $U$ on $\Delta KL$.
Test (10) is a mediator (it explains most of the shape) rather than
a confound (it doesn't make the shape go away). Long-form treatment
of every test is in Appendix~C of the companion paper.}
\label{tab:short-falsification}
\end{table}

The mediator clarifies the substantive interpretation. Filings at
flux extremes are in vocabulary clusters whose peers can also support
real businesses; the flux-neutral middle is in clusters where peers
are also failing. The survival $U$ is a niche-viability signal
mediated through $\Delta KL$. The linear residual --- ``forward flux
still helps after the mediator is included'' --- is the part of
$\Delta KL$'s information that competition-density alone cannot
reproduce.

\section{What we don't know}

(i) Whether the harvest tail reflects intelligent harvesting (firms
recognising an industry transition and extracting from the legacy
side with their eyes open), passive late-occupation (firms whose
business naturally uses the receding vocabulary because they are
still doing the legacy thing), or selection on something unobservable
(firms surviving long enough to file legacy-vocabulary trademarks
have already proved durable). The trademark record does not
distinguish these. One additional confound we cannot fully rule out
in principle is that the harvest tail's category vocabulary recedes
in the corpus precisely because a single dominant firm has captured
the category --- ``Xerox'' becomes the only trademark using
\texttt{xerographic} because everyone else exits the niche, not
because the niche is dying. On the 1990--2021 window this story is
implausible at the scale we observe ($n = 2.6$M harvest-side
debuts across $44$ industries, none of them dominated by a single
firm to that degree), but we flag it explicitly in the long version.
(ii) Why the linear $z(\Delta KL)$ effect on survival
survives the cluster-viability control. The most plausible
explanation is that conditional on a viable cluster, firms whose
vocabulary moves toward the cluster's future are ahead of firms
whose vocabulary moves against it.
(iii) Whether the firm-level financial reward to flux
($+3.2$ pp gross margin per $\sigma$ of trailing $\Delta KL$ for
public firms; $+28$ pp 4-year excess return per $\sigma$ of debut
$\Delta KL$) is robust to the cluster-viability control. We have
not yet computed the firm-level mediation regression.

The companion paper \texttt{paper/main.pdf} reports the same headline
finding with substantially more apparatus: full per-industry detail
(44 NICE classes worked separately), the financial-outcomes panel on
$\sim 1{,}500$ SEC-listed firms, multi-horizon stock returns analysis
with patent-count controls, the cross-buzzword pioneer-ladder
analysis, the canonical-winner-and-failed-firm KL profile, the
per-industry deep-dive appendix with $44$ scatter plots of public
firms by debut year, and full long-form treatment of all ten
falsification tests in Appendix~C.

\end{document}
